% source: https://tex.stackexchange.com/a/766304 (question 472186, score 4, block 1)
\documentclass[border=5pt]{standalone}% compile with lualatex only
\usepackage[svgnames]{xcolor}
\usepackage[3d]{luadraw}%https://github.com/pfradin/luadraw
%https://ctan.org/pkg/luadraw
\usepackage{fourier-otf}
%https://tex.stackexchange.com/questions/472186/how-to-construct-a-circle-that-is-tangent-to-three-given-semi-circles
\begin{document}
\begin{luadraw}{name=tangent_circles}
local ld = luadraw
local cpx = ld.cpx
local Z = cpx.Z
local g = ld.graph:new{ window={-5,5,-4,5.5}, size={10,10} }

local function hyperbola(F1, F2, r, t1, t2) -- returns {M / MF1-MF2 = r}
    local c = cpx.abs(F2-F1)/2
    local O, u = (F1+F2)/2, cpx.normalize(F2-F1)
    local a = r/2
    local b = math.sqrt(c^2-a^2)
    local matrix = {O, u, cpx.I*u}
    local L = ld.parametric(function(t) return Z(a*math.cosh(t), b*math.sinh(t)) end, t1,t2)
    return ld.mtransform(L, matrix)
end

-- data
local A, B, C = Z(-3,-1), Z(2,0), Z(3.5,3) -- Z(2.5,-2), Z(0,4)
-- calculus
local rAB, rBC, rCA = cpx.abs(A-B)/2, cpx.abs(B-C)/2, cpx.abs(C-A)/2
local C1, A1, B1 = (A+B)/2, (B+C)/2, (A+C)/2
local d = 3
local L1 = hyperbola(C1, A1, rBC-rAB,-d,d) 
local L2 = hyperbola(C1, B1, rCA-rAB,-d,d) 
local L3 = hyperbola(A1, B1, rCA-rBC,-d,d) 
local I = ld.interL(L1,L2) -- potential centers
for _, J in ipairs(I) do
    if (cpx.det(J-A,C-A)>=0) and (cpx.det(B-A,J-A)>=0) and (cpx.det(J-C,B-C)>=0)
        then -- J is inside {A,B,C}
            I = J; break 
        end
end
local R = cpx.abs(I-C1)+rAB -- radius
local P1, P2, P3 = I+R*cpx.normalize(A1-I), I+R*cpx.normalize(B1-I), I+R*cpx.normalize(C1-I) -- contact
-- drawing
g:Linewidth(12)
g:Dpolyline({A,B,C}, true, "blue")
g:Darc(A,C1,B,rAB,1, "gray"); g:Darc(B,A1,C,rBC,1, "gray") 
g:Darc(C,B1,A,rCA,1, "gray")
g:Dcircle(I,R,"green")
g:Ddots({P1,P2,P3,I}, "red")
g:Dpolyline( ld.concat(L1,L2,L3), "thin, dashed")
g:Dlabel("$B$",B,{pos="NW"}, "$A$",A,{pos="N",dist=0.1}, "$C$",C,{pos="W"},
    "$P_1$",P1,{pos="E"}, "$P_2$",P2,{pos="N"}, "$P_3$",P3,{pos="S"}, 
    "$I$",I,{pos="S"})
g:Show()
\end{luadraw}
\end{document}
